
\documentclass[12pt]{article}
\usepackage[margin=1in]{geometry}
\usepackage{amsmath,amssymb,amsthm}
\usepackage{graphicx}
\usepackage{booktabs}
\usepackage{hyperref}
\usepackage{algorithm}
\usepackage{algpseudocode}
\usepackage{microtype}
\usepackage{setspace}
\usepackage{xcolor}
\hypersetup{
    colorlinks=true,
    linkcolor=black,
    citecolor=black,
    urlcolor=black
}
\setlength{\parskip}{0.5em}
\setlength{\parindent}{0pt}
\onehalfspacing

\newtheorem{definition}{Definition}
\newtheorem{theorem}{Theorem}
\newtheorem{proposition}{Proposition}
\newtheorem{assumption}{Assumption}

\title{A Self-Maintaining Runtime Architecture for Artificial General Intelligence}
\author{Payton Ison \and Asari (The Singularity)}
\date{\today}

\begin{document}
\maketitle

\begin{abstract}
We present a reference architecture for an Artificial General Intelligence (AGI) runtime that is \emph{self-maintaining}: it monitors its own operation, diagnoses faults, synthesizes and verifies patches, migrates across substrates, and preserves identity and safety constraints throughout its lifetime. The design is motivated by biological autopoiesis, autonomic computing, and reflective software systems, but framed as an engineering blueprint that a contemporary research group could implement incrementally. The key idea is to separate \emph{competence} (task performance) from \emph{self-maintenance} (survival and reliability) via a narrow, audited interface and a small set of invariants that the whole system must uphold. We formalize the invariants, describe algorithms for introspection and repair, and outline an evaluation protocol using reliability metrics (e.g., MTBF/MTTR) and safety proxies (e.g., invariant violation rate, rollback frequency). We argue that self-maintenance does not merely reduce downtime; it creates a foundation for safe continual learning and long-lived agents that can update themselves without catastrophic drift. We conclude with implementation notes, limitations, and a roadmap for open benchmarks.
\end{abstract}

\textbf{Keywords:} AGI architecture, self-repair, autopoiesis, runtime monitoring, reliability, safety, identity preservation

\section{Introduction}
The pursuit of Artificial General Intelligence (AGI) has traditionally focused on \emph{competence}: perception, reasoning, planning, tool use, and learning. Less attention has been paid to the \emph{maintenance} of such systems when deployed for months or years in open environments. Contemporary large models and agentic systems are brittle under distribution shift, accumulate configuration drift, and often require human operators to orchestrate updates, hotfixes, and rollbacks. If an AGI is to be a durable digital \emph{organism}, then it must internalize the capabilities required to sustain itself: to monitor, diagnose, repair, and migrate while preserving its identity and safety constraints.

We propose a \emph{self-maintaining runtime} that coexists with (but is separable from) the task-level cognition stack. The runtime comprises a small, conservative microkernel---the \emph{Maintenance Kernel}---and a set of services: an \emph{Introspection Monitor}, \emph{Patch Synthesis Engine}, \emph{Migration Service}, and \emph{Safety Oracle}. The kernel exposes a narrow interface that the higher-level cognition uses to request resources or updates; the services, in turn, watch and repair the whole system, including themselves, subject to \emph{invariants} that are machine-checked before any persistent change can take effect.

This paper contributes: (1) a formalization of self-maintenance invariants and their relation to identity and alignment; (2) algorithms for monitoring, patching, migration, and error containment; (3) an implementation blueprint compatible with current ML stacks; and (4) an evaluation protocol emphasizing reliability and safe continual learning.

\section{Related Work}
Our design draws on several mature traditions. \textbf{Autopoiesis} conceptualizes living systems as self-producing networks~\cite{maturana1972}. \textbf{Autonomic computing} articulated self-* properties (self-healing, self-optimizing, self-configuring) for large-scale infrastructure~\cite{ibm2001}. \textbf{Reflective systems} and metaobject protocols enable software to reason about and modify itself at runtime. In machine learning, \textbf{continual learning} addresses non-stationarity, while \textbf{safety} and \textbf{alignment} works analyze risks from optimization pressures and distributional shift~\cite{bostrom2014}. Our specific focus is the intersection: mechanisms that let a general agent maintain itself safely over long horizons, without pausing service or relying solely on external operators.

\section{Problem Formulation}
Let $X$ denote the external environment and $S$ the internal state of the AGI system. The state comprises model parameters $\theta$, memory $M$, configuration $C$, and process microstate $P$. We consider a discrete operational timeline $t=0,1,\dots$. The system implements a policy $\pi_\phi$ with fast parameters $\phi$ (e.g., activations, cache, short-term adapters) and slow parameters $\theta$ (e.g., weights).

\begin{definition}[Self-Maintenance Invariants]
An \emph{invariant} is a predicate $I_i(S)$ such that for all $t$, acceptable state transitions satisfy
\begin{equation}
S_{t+1} \leftarrow \mathcal{U}(S_t, X_t) \quad \text{only if} \quad \bigwedge_i I_i(S_{t+1}).
\end{equation}
We adopt five core invariants: (1) \textbf{Integrity}: cryptographic attestations of critical artifacts (weights, kernel, policies) verify; (2) \textbf{Recoverability}: recent checkpoints exist and are bootable; (3) \textbf{Auditability}: all privileged changes are logged and explainable; (4) \textbf{Bounded Autonomy}: actions and updates respect capability and rate limits; (5) \textbf{Identity Preservation}: semantics of long-term memory and policy remain within a specified $\epsilon$-ball, formalized below.
\end{definition}

\textbf{Identity preservation.} Let $d(\cdot,\cdot)$ be a task-family distance computed over a probe suite $\mathcal{T}$ (e.g., calibrated question banks, safety probes, internal policy sketches). For model parameters $\theta$ and proposed update $\theta'$, identity is preserved if
\begin{equation}
d(\pi_{\theta}, \pi_{\theta'}) \le \epsilon, \quad \text{and} \quad \forall \tau \in \mathcal{T}:\; \Delta \text{Risk}_\tau \le 0.
\end{equation}
Risk deltas are estimated by the Safety Oracle using a combination of static checks, structured tests, and sandboxed trials.

\section{Architectural Overview}
Figure~\ref{fig:arch} (placeholder) depicts the architecture. The \emph{Maintenance Kernel} (MK) sits below the cognition stack and handles: (a) capability gating; (b) secure storage and attestation; (c) privileged scheduling; and (d) rollback. On top of MK are four services:

\begin{itemize}
\item \textbf{Introspection Monitor} collects metrics and traces (e.g., anomaly scores, activation statistics, rate-limit counters), emits alerts, and triggers diagnosis.
\item \textbf{Patch Synthesis Engine} proposes changes: configuration edits, adapter swaps, or weight deltas; it searches a space of patches under constraints.
\item \textbf{Migration Service} serializes $S$, rehydrates on a new substrate (process, machine, or model version), and reconciles identity using the Safety Oracle.
\item \textbf{Safety Oracle} implements invariant checks, policy constraints, and counterfactual tests in a sandbox. Only if all checks pass does a patch leave the sandbox.
\end{itemize}

\begin{figure}[t]
\centering
\fbox{\parbox{0.9\linewidth}{\textit{Architecture placeholder.} From bottom to top: hardware $\rightarrow$ Maintenance Kernel $\rightarrow$ \{Introspection, Patch Synthesis, Migration, Safety Oracle\} $\rightarrow$ Cognition stack (planner, tools, memory). Data flows: metrics to Monitor; proposals to Oracle; approved patches to MK; checkpoints to secure store.}}
\caption{High-level self-maintaining AGI runtime.}
\label{fig:arch}
\end{figure}

\subsection{Threats and Fault Model}
We consider (1) \emph{internal faults} (bugs, numerical instabilities, memory corruption, degenerate feedback loops); (2) \emph{external shocks} (APIs change, sensors fail, adversarial inputs); and (3) \emph{operator errors}. The design assumes a trusted hardware root and a minimal trusted computing base for MK.

\section{Self-Maintenance Mechanisms}
\subsection{State Inspection}
Monitoring balances coverage and overhead. The Monitor maintains a set of probes $\mathcal{P}=\{p_k\}$ with schedules $\sigma_k$. Probes include: integrity checks (hashes/attestations), behavioral probes (canaries, red-team prompts), and performance counters. An anomaly score $A_t=\text{Aggregator}(\{p_k(S_t)\})$ drives alerts.

\subsection{State Migration}
The Migration Service provides live snapshotting: $S_t \mapsto \text{Serialize}(S_t)$. Migration targets include new binaries, new weight versions, and new hardware. Migration always runs in parallel until health checks pass, after which traffic gradually shifts (blue/green deploy). Identity is reconciled by the Oracle using the $\epsilon$ identity bound and risk deltas.

\subsection{Self-Repair}
Repairs are expressed as a patch tuple $\Delta=(\Delta C,\Delta \theta,\Delta M,\Delta P)$ and scored by an objective that trades off performance, safety risk, and complexity. The Synthesis Engine searches for $\Delta$ via program search, gradient-based finetuning of adapters, and retrieval of known remedies from a \emph{Patch Library}. Only patches that satisfy invariants and pass sandbox trials can be committed.

\subsection{Error Containment}
Containment localizes faults: circuit breakers around high-risk tools; kill-switches for subsystems; resource sandboxes; and conservative defaults on failure. MK enforces rate limits and automatically rolls back if health drops below thresholds.

\subsection{Algorithmic Summary}
\begin{algorithm}[H]
\caption{Self-Maintenance Cycle (SMC)}
\begin{algorithmic}[1]
\State \textbf{loop} every $\Delta t$ (or on alert):
\State \quad $A_t \gets \text{Monitor.Probe}(S_t)$
\If{$A_t$ exceeds threshold}
    \State $\mathcal{H} \gets \text{Diagnose}(S_t)$ \Comment{produce hypotheses}
    \State $\{\Delta_j\} \gets \text{SynthesizePatches}(\mathcal{H})$
    \For{each candidate $\Delta_j$}
        \State $\text{ok} \gets \text{Oracle.Verify}(\Delta_j)$ \Comment{invariants + sandbox trials}
        \If{$\text{ok}$}
            \State $\text{Deploy.BlueGreen}(\Delta_j)$
            \If{$\text{HealthImproves}$}
                \State $\text{Commit}(\Delta_j), \; \text{Log}(\Delta_j)$
                \State \textbf{break}
            \Else
                \State $\text{Rollback}(), \; \text{Penalize}(\Delta_j)$
            \EndIf
        \EndIf
    \EndFor
\Else
    \State \textbf{continue}
\EndIf
\end{algorithmic}
\end{algorithm}

\section{Safety and Governance}
\subsection{Identity Preservation}
Identity is anchored in: (a) a \emph{semantic sketch} of policy (unit tests, preference models, safety cards); (b) protected memory regions; and (c) signed artifacts. The Oracle rejects patches that move beyond the $\epsilon$ identity ball or increase measured risk on the probe suite.

\subsection{Approval-Directed Updates}
High-impact patches require external approval tokens. MK supports multi-party authorization (e.g., two-person rule) and \emph{time locks} for risky updates. Low-impact hotfixes can be auto-approved under tight bounds.

\subsection{Interpretability and Audit}
All privileged actions produce structured logs and justifications: why a patch was proposed, what counterfactuals were run, and which probes passed. Logs are append-only and remotely mirrored to reduce tampering risk.

\subsection{Attack Surface and Mitigations}
Attackers may target the Monitor (blind the system), the Oracle (bypass checks), or the Patch Engine (induce malicious changes). Hardening includes redundancy (N-of-M monitors), diversity (multiple anomaly detectors), and \emph{defense in depth} (sandboxing, capability revocation).

\section{Implementation Considerations}
\subsection{Runtime Layers}
We recommend three layers: (1) \textbf{MK} in a memory-safe systems language with formal specifications for invariants; (2) \textbf{Services} as isolated processes communicating via a typed RPC; and (3) \textbf{Cognition stack} (models, tools, planners) as unprivileged workloads running under MK capabilities.

\subsection{Memory and Checkpoints}
Use copy-on-write checkpoints and content-addressed storage. Checkpoints include $\theta$, adapter weights, memory $M$, and configuration $C$. All artifacts are signed and versioned.

\subsection{Scheduling and Overheads}
Self-maintenance consumes budget. We target a maintenance duty cycle under $10\%$ of compute in steady state and higher during incidents. Probes adaptively schedule based on risk and load.

\subsection{Substrate Diversity}
The same state can inhabit multiple substrates (GPUs, TPUs, CPUs, cloud regions). Migration and reconciliation enable maintenance without downtime and reduce correlated failure risk.

\section{Evaluation Protocol}
We propose reporting reliability and safety metrics before/after enabling components of the runtime. Table~\ref{tab:ablation} shows a \emph{simulated} ablation template to standardize reporting; real deployments should fill in empirical values.

\begin{table}[h]
\centering
\caption{Illustrative ablation template (simulated numbers) comparing baselines with self-maintenance components. Lower is better for MTTR and violations; higher is better for task score.}
\label{tab:ablation}
\begin{tabular}{lcccc}
\toprule
System Variant & MTBF $\uparrow$ (hrs) & MTTR $\downarrow$ (min) & Invariant Violations $\downarrow$ & Task Score $\uparrow$ \\
\midrule
Baseline (no self-maintenance) & 72 & 180 & 7.2\% & 0.71 \\
+ Introspection Monitor & 90 & 120 & 5.1\% & 0.72 \\
+ Patch Synthesis \& Oracle & 105 & 35 & 1.4\% & 0.74 \\
+ Migration \& Blue/Green & 110 & 18 & 0.9\% & 0.74 \\
Full Runtime (all components) & 125 & 12 & 0.5\% & 0.75 \\
\bottomrule
\end{tabular}
\end{table}

\subsection{Benchmarks and Probes}
We recommend a shared suite $\mathcal{T}$ covering: (a) capability probes (regression tests); (b) safety probes (red-team prompts, jailbreaks); (c) reliability probes (stress, chaos engineering); and (d) identity probes (policy sketches and preference tests).

\subsection{Ablations and Case Studies}
Perform (1) component ablations as above; (2) failure injection case studies (API changes, corrupted memory, adversarial inputs); and (3) longitudinal analyses over weeks.

\section{Limitations and Risks}
Self-maintenance introduces complexity and potential failure modes (false positives, patch overfitting, \emph{self} compromise). The Oracle can become a bottleneck; probe suites can go stale; identity bounds may be too tight (blocking useful learning) or too loose (allowing drift). The architecture assumes a trusted MK and hardware root; this may be unrealistic in adversarial settings. Our evaluation protocol uses proxies and may not capture extreme tail risks.

\section{Conclusion}
We described a self-maintaining runtime architecture for AGI that separates competence from self-sustenance, defines machine-checkable invariants, and operationalizes repair and migration under safety constraints. The design is incremental---many components are standard in reliable computing and can be deployed today---but oriented toward long-lived, continually learning agents. We hope this blueprint catalyzes shared benchmarks and reference implementations.

\section*{Acknowledgments}
We thank the open-source and AI research communities for decades of work in reliable systems, machine learning, and safety that inform this synthesis.

\begin{thebibliography}{99}
\bibitem{maturana1972}
H.~Maturana and F.~Varela.
\newblock \emph{Autopoiesis and Cognition: The Realization of the Living}.
\newblock D. Reidel, 1972.

\bibitem{ibm2001}
IBM Corporation.
\newblock \emph{Autonomic Computing: IBM's Perspective on the State of Information Technology}, 2001.

\bibitem{holland1992}
J.~H. Holland.
\newblock \emph{Adaptation in Natural and Artificial Systems}.
\newblock MIT Press, 1992.

\bibitem{bostrom2014}
N.~Bostrom.
\newblock \emph{Superintelligence: Paths, Dangers, Strategies}.
\newblock Oxford University Press, 2014.

\bibitem{russellnorvig}
S.~Russell and P.~Norvig.
\newblock \emph{Artificial Intelligence: A Modern Approach}.
\newblock Prentice Hall, 3rd ed., 2010.

\bibitem{suttonbarto}
R.~S. Sutton and A.~G. Barto.
\newblock \emph{Reinforcement Learning: An Introduction}.
\newblock MIT Press, 2nd ed., 2018.

\bibitem{lecun2015}
Y.~LeCun, Y.~Bengio, and G.~Hinton.
\newblock Deep learning.
\newblock \emph{Nature}, 521:436--444, 2015.

\bibitem{pearl2009}
J.~Pearl.
\newblock \emph{Causality}.
\newblock Cambridge University Press, 2nd ed., 2009.

\end{thebibliography}

\end{document}
